% GENERATED FILE. Edit canonical lessons, then run npm run generate:books.
% canonical-source-hash: fnv1a64:c326f1c03bbdafa4
% canonical-lessons: GE-C344-beeinflussen, GE-C344-einfluss, GE-C344-auswirkung, GE-C344-auswirken, GE-C344-betreffen

\chapter{Einfluss haben}
\label{ch:ge-einfluss-haben}
\hlchaptermodality{344}

\begin{chapteropening}
\textbf{By the end of this chapter:} \emph{I can say what influences something, what impact it has and whom it concerns.}

Complete the last lesson of chapter 344: I can say what influences something, what impact it has and whom it concerns.
\end{chapteropening}

\section[beeinflussen]{beeinflussen — to influence}

\label{lesson:GE-C344-beeinflussen}



\begin{quote}
Before the new one: say the German for so that, then the German for therefore.
\end{quote}

\subsection*{You'll want to know: beeinflussen}
\textbf{beeinflussen} — \textquotedblleft{}to influence\textquotedblright{}.

\begin{quote}
Das Wetter beeinflusst meine Laune.
\end{quote}

\begin{quote}
\emph{The weather influences my mood.}
\end{quote}

Inseparable: \emph{es beeinflusst, es hat beeinflusst}. It takes a direct object.

\begin{cousinweb}
\emph{beeinflussen} is built on \emph{Einfluss} \textquotedblleft{}influence\textquotedblright{}, literally an \textquotedblleft{}inflow\textquotedblright{}, from \emph{ein} \textquotedblleft{}in\textquotedblright{} and \emph{Fluss} \textquotedblleft{}flow, river\textquotedblright{}. \emph{Einfluss} is usually explained as a rendering of Latin \emph{influxus} \textquotedblleft{}a flowing in\textquotedblright{}; EtymDB has no history for it.
\end{cousinweb}

\begin{grammarlens}[title={What you've built}]
One more word for telling, explaining and describing at B1.
\end{grammarlens}

\subsection*{Your turn}
\begin{itemize}
\raggedright
  \item \emph{Say it:} \emph{beeinflussen}
  \item \emph{Say it:} \emph{beeinflussen}, once more
  \item \emph{Say it:} a sentence of your own with \emph{beeinflussen}
  \item \emph{Recall:} say the German for so that, then the German for therefore, then say \emph{beeinflussen} again
\end{itemize}

\begin{tcolorbox}[breakable,colback=teal!4,colframe=teal!35!black,title={Before you move on}]
What does beeinflussen mean? (\textquotedblleft{}To influence\textquotedblright{}.) Say it once more.
\end{tcolorbox}
\section[der Einfluss]{der Einfluss — the influence}

\label{lesson:GE-C344-einfluss}



\begin{quote}
Before the new one: say the German for therefore, then the German for to influence.
\end{quote}

\subsection*{You'll want to know: der Einfluss}
\textbf{der Einfluss} — \textquotedblleft{}the influence\textquotedblright{}.

\begin{quote}
Meine Großmutter hatte einen großen Einfluss auf mich.
\end{quote}

\begin{quote}
\emph{My grandmother had a great influence on me.}
\end{quote}

Plural: \emph{die Einflüsse}. \emph{Einfluss auf} + accusative.

\begin{cousinweb}
\emph{Einfluss} joins \emph{ein} \textquotedblleft{}in\textquotedblright{} to \emph{Fluss} \textquotedblleft{}flow, river\textquotedblright{}. It is usually explained as a rendering of Latin \emph{influxus} \textquotedblleft{}a flowing in\textquotedblright{}, which late Latin also used of the influence of the stars (Lewis \& Short); EtymDB has no history for it.
\end{cousinweb}

\begin{grammarlens}[title={What you've built}]
One more word for telling, explaining and describing at B1.
\end{grammarlens}

\subsection*{Your turn}
\begin{itemize}
\raggedright
  \item \emph{Say it:} \emph{der Einfluss}
  \item \emph{Say it:} \emph{der Einfluss}, once more
  \item \emph{Say it:} a sentence of your own with \emph{der Einfluss}
  \item \emph{Recall:} say the German for therefore, then the German for to influence, then say \emph{der Einfluss} again
\end{itemize}

\begin{tcolorbox}[breakable,colback=teal!4,colframe=teal!35!black,title={Before you move on}]
What does der Einfluss mean? (\textquotedblleft{}The influence\textquotedblright{}.) Say it once more.
\end{tcolorbox}
\section[die Auswirkung]{die Auswirkung — the impact, the effect (on something)}

\label{lesson:GE-C344-auswirkung}



\begin{quote}
Before the new one: say the German for to influence, then the German for the influence.
\end{quote}

\subsection*{You'll want to know: die Auswirkung}
\textbf{die Auswirkung} — \textquotedblleft{}the impact, the effect (on something)\textquotedblright{}.

\begin{quote}
Die neue Straße hat Auswirkungen auf das ganze Dorf.
\end{quote}

\begin{quote}
\emph{The new road has an impact on the whole village.}
\end{quote}

Mostly plural: \emph{die Auswirkungen auf} + accusative.

\begin{cousinweb}
\emph{Auswirkung} is the noun of \emph{sich auswirken}, built on \emph{wirken} \textquotedblleft{}to work\textquotedblright{}, inherited from Old High German \emph{wirken} and Proto-Germanic \emph{wirkijaną} (EtymDB).
\end{cousinweb}

\begin{grammarlens}[title={What you've built}]
One more word for telling, explaining and describing at B1.
\end{grammarlens}

\subsection*{Your turn}
\begin{itemize}
\raggedright
  \item \emph{Say it:} \emph{die Auswirkung}
  \item \emph{Say it:} \emph{die Auswirkung}, once more
  \item \emph{Say it:} a sentence of your own with \emph{die Auswirkung}
  \item \emph{Recall:} say the German for to influence, then the German for the influence, then say \emph{die Auswirkung} again
\end{itemize}

\begin{tcolorbox}[breakable,colback=teal!4,colframe=teal!35!black,title={Before you move on}]
What does die Auswirkung mean? (\textquotedblleft{}The impact, the effect (on something)\textquotedblright{}.) Say it once more.
\end{tcolorbox}
\section[auswirken]{auswirken — to affect (reflexive, sich auswirken auf)}

\label{lesson:GE-C344-auswirken}



\begin{quote}
Before the new one: say the German for the influence, then the German for the impact.
\end{quote}

\subsection*{You'll want to know: auswirken}
\textbf{auswirken} — \textquotedblleft{}to affect (reflexive, sich auswirken auf)\textquotedblright{}.

\begin{quote}
Der Stress wirkt sich auf meinen Schlaf aus.
\end{quote}

\begin{quote}
\emph{The stress affects my sleep.}
\end{quote}

Reflexive and separable: \emph{Das wirkt sich auf … aus}, \emph{das hat sich ausgewirkt}.

\begin{cousinweb}
\emph{auswirken} is built from \emph{aus} \textquotedblleft{}out\textquotedblright{} and \emph{wirken} \textquotedblleft{}to work\textquotedblright{}, inherited from Old High German \emph{wirken} and Proto-Germanic \emph{wirkijaną} \textquotedblleft{}to work, to make\textquotedblright{} (EtymDB): an effect works its way out.
\end{cousinweb}

\begin{grammarlens}[title={What you've built}]
One more word for telling, explaining and describing at B1.
\end{grammarlens}

\subsection*{Your turn}
\begin{itemize}
\raggedright
  \item \emph{Say it:} \emph{auswirken}
  \item \emph{Say it:} \emph{auswirken}, once more
  \item \emph{Say it:} a sentence of your own with \emph{auswirken}
  \item \emph{Recall:} say the German for the influence, then the German for the impact, then say \emph{auswirken} again
\end{itemize}

\begin{tcolorbox}[breakable,colback=teal!4,colframe=teal!35!black,title={Before you move on}]
What does auswirken mean? (\textquotedblleft{}To affect (reflexive, sich auswirken auf)\textquotedblright{}.) Say it once more.
\end{tcolorbox}
\section[betreffen]{betreffen — to concern, to affect}

\label{lesson:GE-C344-betreffen}



\begin{quote}
Before the new one: say the German for the impact, then the German for to affect.
\end{quote}

\subsection*{You'll want to know: betreffen}
\textbf{betreffen} — \textquotedblleft{}to concern, to affect\textquotedblright{}.

\begin{quote}
Die neue Regel betrifft alle Schüler.
\end{quote}

\begin{quote}
\emph{The new rule concerns all pupils.}
\end{quote}

Inseparable and strong: \emph{es betrifft, es hat betroffen}. \emph{Was mich betrifft, …} — as far as I am concerned, …

\begin{cousinweb}
\emph{betreffen} is built with \emph{be-} on \emph{treffen} \textquotedblleft{}to hit, to meet\textquotedblright{}, inherited from Old High German \emph{treffan} \textquotedblleft{}to hit\textquotedblright{} and Proto-Germanic \emph{drepaną} (EtymDB): a rule hits the people it is about.
\end{cousinweb}

\begin{grammarlens}[title={What you've built}]
One more word for telling, explaining and describing at B1.
\end{grammarlens}

\subsection*{Your turn}
\begin{itemize}
\raggedright
  \item \emph{Say it:} \emph{betreffen}
  \item \emph{Say it:} \emph{betreffen}, once more
  \item \emph{Say it:} a sentence of your own with \emph{betreffen}
  \item \emph{Recall:} say the German for the impact, then the German for to affect, then say \emph{betreffen} again
\end{itemize}

\begin{tcolorbox}[breakable,colback=teal!4,colframe=teal!35!black,title={Before you move on}]
What does betreffen mean? (\textquotedblleft{}To concern, to affect\textquotedblright{}.) Say it once more.
\end{tcolorbox}
